\chapter{fast vs Runner：为什么两个都要懂}

\section{本章要解决的困惑}

fast 和 Runner 都能“回测”，但它们不是同一种工具。

\begin{longtable}{p{0.20\textwidth}p{0.34\textwidth}p{0.36\textwidth}}
\toprule
工具 & 回答什么 & 不回答什么 \\
\midrule
fast & 状态机在历史 decision frame 上表现如何 & 真实订单到达、TCP、完整 event log \\
Runner & 本地交易所如何处理 intent、latency、fill、portfolio & 大批量扫因子和 profile \\
\bottomrule
\end{longtable}

\section{fast 适合什么}

fast 适合：

\begin{itemize}[leftmargin=2em]
  \item 看 entry trigger 是否有基本 edge；
  \item 比较 spread q70 gate；
  \item 比较 gamma/capacity profile；
  \item 看 daily stability；
  \item 快速检查 mid vs executable。
\end{itemize}

它不适合证明 live path 一定正确。

\section{Runner 适合什么}

Runner 适合：

\begin{itemize}[leftmargin=2em]
  \item 验证 Bot 与 Runner event contract；
  \item 验证 latency 和 arrival quote；
  \item 验证 fill price；
  \item 生成可审计 event log；
  \item 检查 Monitor 的因果链展示。
\end{itemize}

它不适合扫大量因子，因为慢且工程噪声更多。

\section{两个结果怎么比较}

比较前必须确认：

\begin{lstlisting}
same symbol
same date range
same decision clock
same policy
same admission threshold
same capacity profile
same exit/fill model
\end{lstlisting}

如果 fast 是 \code{fixed60_mid}，Runner 是 \code{top_of_book_taker_ioc_v1}，净收益不同不一定是 bug。

\section{什么时候先 fast，什么时候上 Runner}

先 fast：

\begin{itemize}[leftmargin=2em]
  \item 新 entry 条件；
  \item 新 gamma profile；
  \item 新 admission threshold；
  \item 只看研究方向是否值得继续。
\end{itemize}

上 Runner：

\begin{itemize}[leftmargin=2em]
  \item 改了 Bot message loop；
  \item 改了 Runner fill/latency/order code；
  \item 要验收 Monitor；
  \item 要证明某个 profile 能在 strict path 下复现。
\end{itemize}

\checkpoint{fast 是研究显微镜，Runner 是本地交易所。把它们混成一个东西，文档和代码都会变难。}

